\section{The remaining unitary groups}
\label{sec:unitary-induction}

We now pass from the odd-degree division case to arbitrary outer forms of
type~$A$.  The induction is on the \emph{reduced degree}, simultaneously
over all number fields.  Thus a special unitary group of an algebra of
degree~$m$ over a finite extension of~$k$ is an admissible induction input
whenever $m$ is smaller than the degree currently under consideration.
Restriction of scalars does not change this convention.

\begin{theorem}\label{thm:outer-a}
The Margulis--Platonov assertion holds for every simply connected
absolutely almost simple group of outer type~$A$ over a number field.
\end{theorem}

Write $L/k$ for a quadratic extension, $D$ for a central simple
$L$-algebra of degree~$n$, and $*$ for a unitary involution on~$D$.
Set $S=\SU(D,*)$ and $U=\U(D,*)$.  We continue to write $i\in L^\times$
for a fixed element with $i^*=-i$.  By the established isotropic cases
and Proposition~\ref{prop:finite-defect}, it suffices to assume that $S$
is $k$-anisotropic and prove $V_0/R=1$, where $R=S(k)^e$ is an abstract
power subgroup.  In this section a group element is said to be
\emph{killed} when its image in the finite group $S(k)/R$ is trivial.
Lemma~\ref{lem:restriction-mp} will be used for smaller special groups
and their restrictions of scalars.  Its local power criterion applies
to their images in~$S$ as well.

The degree-two case is an established inner type~$A_1$ case of
Theorem~\ref{thm:known-mp}.  For odd division degree,
Proposition~\ref{prop:color-dichotomy} and
Theorems~\ref{thm:no-characters} and~\ref{thm:no-simple-colors} give
$V_0/R=1$.  Three arguments remain: a ternary split-algebra case, the
odd matrix-algebra case, and the even-degree induction step.

\subsection{Uniform local power roots}

Some of the smaller subgroups below vary with the element being
factored.  The following observation allows their local neighborhoods
to be chosen in the fixed ambient algebra before those subgroups are
constructed.

\begin{lemma}\label{lem:unitary-uniform-roots}
Fix a completion $k_v$, a finite-dimensional matrix realization of~$D$,
and an integer $e\geq1$.  There is a neighborhood of~$1$ with the
following property.  An element in that neighborhood which belongs to
one of the unitary or special unitary subgroups arising from a
$*$-stable subalgebra of~$D$, or from an invariant summand of its
Hermitian module, has an $e$-th root in the same subgroup.  The
neighborhood can be chosen uniformly for all such subalgebras and
summands of the bounded degrees occurring here.
\end{lemma}

\begin{proof}
In a sufficiently small matrix neighborhood, the convergent series
$(1+X)^{1/e}$ defines the root close to~$1$.  A finite-dimensional local
subalgebra is closed, so the series stays in every subalgebra containing
$X$.  Uniqueness of the root close to~$1$ gives
$(a^{1/e})^*=(a^*)^{1/e}=(a^{1/e})^{-1}$ when $a^*=a^{-1}$.
The same series is the identity on a summand where $a$ is the identity.

For a special group, the reduced determinant of the root has $e$-th
power~$1$.  All its reduced eigenvalues occur among the ambient
eigenvalues, and all are close to~$1$ over the local algebraic closure
when the ambient matrix is sufficiently close to~$1$.  There are at
most $n$ eigenvalues in each reduced determinant.  Shrinking the fixed
neighborhood therefore makes all the relevant determinants lie in a
neighborhood containing no $e$-th root of unity other than~$1$.  This
proves the special condition uniformly, including after restriction of
scalars.
\end{proof}

\subsection{Ternary Hermitian forms}

\begin{proposition}\label{prop:unitary-ternary}
If $D=M_3(L)$, then $V_0/R=1$.
\end{proposition}

\begin{proof}
At every finite place a ternary Hermitian space is isotropic, so $A$ is
empty and $V_0=S(k)$.  We may assume that the global form is anisotropic.
Use the convention that its pairing $\langle\ ,\ \rangle$ is linear in
the first variable.  Scale the form so that a fixed vector $x$ has norm
$1$, and choose $e_2\perp x$ with norm $\kappa\ne0$.  Let $\Delta$ be
the determinant of the full form in fixed volume coordinates.

Choose a finite set $B$ containing the archimedean places, all places
where $-\Delta$ or $-\kappa$ is not a local norm from~$L$, and the finite
rank-zero places of the special stabilizer of~$x$.  The first two norm
conditions exclude only finitely many finite places: outside the
ramification and the divisors of the constants, units are norms from
an unramified quadratic extension.  Enlarge $B$ for fixed models as
necessary.  We shall factor a representative of a prescribed coset
into three binary special transformations, all locally small at~$B$.

For a representative $u$, put
\[
 z=ux,\qquad \alpha=\langle x,z\rangle,\qquad
 \delta=1-N_{L/k}(\alpha),\qquad
 k_0=2-\alpha-\bar\alpha.
\]
We require $\delta k_0(\delta-k_0)\ne0$.  In the nondegenerate plane
spanned by $x,z$, there is a unique vector $y_*$ whose pairing with each
of $x,z$ is~$1$; inversion of its two-by-two Gram matrix gives
\[
 \langle y_*,y_*\rangle=k_0/\delta.
\]
Choose the vector $v\perp x,z$ with $\det(x,z,v)=1$.  The determinant
of the Gram matrix in this basis gives
$\langle v,v\rangle=\Delta/\delta$.  Consequently a vector
$y=a y_*+b v$ has norm~$1$, has both marked pairings equal to~$a$, and
satisfies $1-N_{L/k}(a)=\kappa N_{L/k}(c)$ precisely when
\begin{equation}\label{eq:unitary-three-norms}
 t=N_{L/k}(a),\qquad
 \frac{\delta-k_0t}{\Delta}=N_{L/k}(b),\qquad
 \frac{1-t}{\kappa}=N_{L/k}(c).
\end{equation}
We insist that $a,b,c$ be nonzero.

We first produce local models at the places in~$B$.  Choose
$y=a x+c e_2$ close to~$x$, with $a,c\ne0$ and norm~$1$; the local norm
map is a submersion at~$1$, so this is possible with $a$ close to~$1$
and $c$ small.  The special isometry $w_1$ of the plane $\langle x,y\rangle$
sending $x$ to~$y$, extended by the identity, is close to~$1$.
Choose $\rho$ close to~$1$ in the special stabilizer of~$y$, and set
\[
 z=\rho x,\qquad
 w_2=\rho w_1^{-1}\rho^{-1},\qquad u_0=w_2w_1.
\]
Then $w_2y=z$, $u_0x=z$, and
$\langle y,x\rangle=\langle y,z\rangle=a$.
We can also arrange linear independence of $x,y,z$ and the three
inequalities imposed above.  Indeed the projection of $x$ onto
$y^\perp$ is nonisotropic, and its generic image under the special
stabilizer gives linear independence.  The scalar inequalities are
nonempty algebraic conditions: taking minus the identity on
$y^\perp$ gives $\alpha=2N_{L/k}(a)-1$, and a generic choice of~$a$
avoids their zeros.  Local points near~$1$ are Zariski dense, so these
nonempty algebraic openings can be met while retaining all the
smallness conditions.

At these models Equation~\eqref{eq:unitary-three-norms} has all three
norm arguments nonzero.  Since a quadratic norm map on nonzero
elements is a local submersion, solutions persist for nearby $u$, even
with $t$ fixed.  Proposition~\ref{prop:finite-defect} now supplies a
representative of the prescribed coset sufficiently close to these
models at~$B$.  At a place outside~$B$, choose $c\ne0$ small and put
$t=1-\kappa N_{L/k}(c)$.  The first factor is then close to~$1$; the
second is close to a nonzero norm because
\[
 \delta-k_0=-N_{L/k}(1-\alpha)
 \quad\hbox{and}\quad -\Delta\in N_{L/k}(L_v^\times).
\]
Thus Equation~\eqref{eq:unitary-three-norms} has the required local
solutions everywhere, with the prescribed ones at~$B$.

Here is the global step, including the passage from a product of
norms to the three separate norms.  Write
\[
 f_1(t)=t,\qquad f_2(t)=(\delta-k_0t)/\Delta,\qquad
 f_3(t)=(1-t)/\kappa.
\]
Their roots $0,\delta/k_0,1$ are distinct.  A smooth projective model
of
\begin{equation}\label{eq:unitary-chatelet}
 N_{L/k}(\xi)=f_1(t)f_2(t)f_3(t)
\end{equation}
is a Ch\^atelet surface.  The theorem of
Colliot-Th\'el\`ene--Sansuc--Swinnerton-Dyer identifies the closure of
its rational points with its Brauer--Manin set
\cite{CTSansucSD}; see also \cite[Section~5.2]{CT1992}.
For a split cubic, its
Brauer group modulo constants is generated by the quaternion symbols
of the quadratic extension and the linear factors
\cite{CTSansucSD}.  The local points just constructed are orthogonal
to these classes: each $f_j(t)$ is individually a norm, so a symbol
written with $f_j$ evaluates to zero.  Replacing $f_j$ by its monic
linear factor only adds a constant Brauer class, and the sum of the
local invariants of that constant is zero by reciprocity.  These
choices give an adelic point; integral choices are available at all
but finitely many places.

Enlarge the finite approximation set to contain $B$, the ramified
places of $L/k$, and all divisors of the leading coefficients, nonzero
roots, and pairwise root differences of the three linear factors.  Apply the
Ch\^atelet theorem to approximate the chosen local points there, on
the affine open where $f_1f_2f_3\ne0$.  At a remaining nonsplit place
the extension is unramified and all these nonzero constants are units.  If
$t$ is integral, at most one $f_j(t)$ has positive valuation.  If $t$
is not integral, all three $f_j(t)$ have valuation $\ord_v(t)$.  Since
their product is a norm by Equation~\eqref{eq:unitary-chatelet}, the
nonzero valuation in the first case, or the common valuation in the
second case, is even.  Units are norms in an unramified quadratic
extension.  Thus every $f_j(t)$ is separately a local norm at every
place, including the approximated places.  The quadratic Hasse norm
theorem \cite[Chapter~VIII, Theorem~3.1]{MilneCFT} gives global $a,b,c$ in
Equation~\eqref{eq:unitary-three-norms}.  Its norm torsors are rational
conics once soluble, so weak approximation gives the prescribed
approximations to $a,b,c$ at~$B$.

For the resulting $y$, both planes $\langle x,y\rangle$ and
$\langle z,y\rangle$ have Gram determinant
$1-N_{L/k}(a)=\kappa N_{L/k}(c)\ne0$.  Their special groups have no
finite rank-zero place outside~$B$, since $-\kappa$ is a local norm
there.  In a nondegenerate binary Hermitian plane, an isometry between
marked norm-one vectors has a unique special extension: extend it to
the orthogonal line and adjust the determinant on that line.  In
local orthogonal coordinates this extension varies smoothly with the
plane and marked vectors.  Hence the two resulting transformations
$w_1,w_2$ approximate the prescribed local transformations, and the
remaining factor $(w_2w_1)^{-1}u$ fixes~$x$ and is close to~$1$ at~$B$.
Lemma~\ref{lem:unitary-uniform-roots}, followed by the established
$A_1$ case and Lemma~\ref{lem:restriction-mp}, kills all three factors.
This kills the prescribed coset.
\end{proof}

\subsection{Odd matrix degree}

Reduction to Hermitian dimension two is an established strategy in the
unitary problem; see Tomanov \cite[p.~895]{Tomanov1992}.
Here the induction parameter is the reduced degree over all number fields.
Suppose $n$ is odd and $D=M_l(\Delta_0)$, where $\Delta_0$ is a division
algebra of degree $d$ and $l>1$.  Then $l\geq3$ and $n=ld$.  The
classification of unitary involutions realizes $S$ as the special
group of a Hermitian $l$-space over a unitary division algebra
\cite[Theorem~4.2(2)]{KMRT}.  Anisotropy implies that every subspace is nondegenerate.

The group is generated by its special transformations on
two-dimensional $\Delta_0$-subspaces.  To see this, choose a nonzero
vector $x$ and an isometry~$u$.  In a plane containing $x,ux$, Witt's
extension theorem gives an isometry sending $x$ to~$ux$.  Its
determinant can be adjusted on the orthogonal line in that plane.
Here the unitary group of a nondegenerate $\Delta_0$-line maps onto
$L^1$ by reduced determinant.  At a nonsplit finite place, the
existence of a unitary involution makes the Brauer corestriction of
the algebra zero.  For a quadratic extension of nonarchimedean local
fields, corestriction preserves the local Brauer invariant, so the
algebra itself is split.  The line therefore becomes an ordinary
Hermitian $d$-space, and its unitary determinant is onto by varying
one orthogonal coordinate.  At a nonsplit real place the extension
is $\bbC/\bbR$, the algebra is again split, and the same argument
works for every Hermitian signature.  At split finite places use
reduced norm surjectivity.  At split real places the division degree is
odd, so there is no quaternionic sign restriction.  Globally, for
$d>1$, each determinant fiber is a torsor under the simply connected
special group of that line; Theorem~\ref{thm:foundations-hasse} gives
its rational point.  For $d=1$ surjectivity is immediate.  Dividing by the resulting
special plane move fixes~$x$, and repetition in its orthogonal
complement proves the generation assertion.

If $d>1$, these plane groups have reduced degree $2d<n$.  They have no
finite rank-zero places.  At a split place this follows from their
two-dimensional module over the local division algebra; at a
nonsplit finite place the algebra is split and the Hermitian
dimension $2d\geq6$ is isotropic.  The induction hypothesis therefore
kills their rational points.  If $d=1$, put each plane move inside a
three-dimensional subspace and apply
Proposition~\ref{prop:unitary-ternary}.  This proves the odd
matrix-algebra case.

\subsection{A quadratic Hermitian subfield in even degree}

Assume now that $n=2m$ with $m\geq2$.  We first construct a fixed
quadratic subfield $S_0\subset D$, disjoint from~$L$, on which $*$ is
the identity.  The centralizer of~$S_0$ will supply the smaller special
unitary group in the induction.

\begin{lemma}\label{lem:unitary-hermitian-quadratic}
There is a quadratic field $S_0=k(s)\subset D$, disjoint from~$L$,
with $s^*=s$ and $s^2=r\in k^\times$, such that the two geometric
eigenspaces of~$s$ have dimension~$m$.
\end{lemma}

\begin{proof}
Choose $r\in N_{L/k}(L^\times)$ with the following local conditions.
At the finitely many exceptional nonsplit places, require $r$ to be a
square.  At a place split in~$L$ where the local index of~$D$ does not
divide~$m$, require $r$ to be a nonsquare, including the corresponding
real sign condition when needed.  At one additional finite place
split in~$L$, require $r$ to be a nonsquare.  This last condition
ensures that $S_0$ is a field and differs from~$L$.  The requirements
are compatible: at split places the norm map is surjective, and a
square can be approximated by norms at nonsplit places.  Weak
approximation in~$L$, followed by the norm, supplies~$r$.

We verify local embeddings of $k_v[s]/(s^2-r)$ with the stated
multiplicities and involution.  At split places, a quadratic field
extension reduces the local division index by its factor of two,
whenever this is required to make the index divide~$m$.  This is
exactly the local embedding criterion for a degree-two etale algebra
in a central simple algebra of degree~$2m$.  In the split quadratic
case one uses two blocks of degree~$m$.  The two components over~$L$
are paired by the involution.

At a nonsplit place where $S_0$ splits, decompose the Hermitian space
as an orthogonal sum of two $m$-spaces.  At every other nonsplit place
we may, by the choice of the exceptional set, assume that $L/k$ is
unramified and the Hermitian form is unimodular.  Its determinant is
a norm.  If $S_0$ locally equals~$L$, the form is hyperbolic and dual
isotropic $m$-spaces give the embedding.  If the two quadratic fields
are distinct, use the orthogonal sum of $m$ copies of one half the
trace to~$L$ of the standard Hermitian line over $LS_0/S_0$.  In the
basis $1,s$ its determinant is~$r$, a norm, so the local
classification of Hermitian forms identifies this form with the
given one.  These facts about algebras and local Hermitian forms are
the usual local classification and embedding theorems
\cite{Reiner,KMRT}.

Let $X$ be the variety of such embeddings, equivalently the images
of~$s$ with $s^*=s$, $s^2=r$, and the fixed geometric multiplicities.
Over the algebraic closure it is a single orbit under~$S$, with
connected stabilizer
\[
 \mathrm{S}(\GL_m\times\GL_m).
\]
The units on $\overline X$ are constant, and, because the ambient
group is simply connected, $\Pic(\overline X)$ is the character
lattice of this stabilizer, of rank one
\cite[Proposition~6.10, equation~(6.10.1)]{Sansuc}. The Hochschild--Serre edge
map gives an injection
$\Br_a X\longrightarrow H^1(k,\Pic(\overline X))$, where
$\Br_a X=\Br_1X/\Br_0X$ and $\Br_0X=\im(\Br k\to\Br X)$;
its kernel before taking this quotient
is precisely the image of the constant classes.  This map commutes
with localization. Define
\[
 B(X)=\ker\!\left(\Br_a X\longrightarrow
                       \prod_v\Br_a X_{k_v}\right).
\]
These are the algebraic Brauer classes locally trivial modulo constants.
Hence $B(X)$ injects into
$\Sha^1(k,\Pic(\overline X))$.  This last group vanishes.  Indeed a
rank-one lattice has trivial or sign Galois action.  The trivial case
has $H^1(k,\bbZ)=0$; in the sign case any nonzero class factors
through the quadratic sign quotient and is detected at an inert
unramified place by Chebotarev.  Inflation adds no classes since a
continuous homomorphism from a profinite group to~$\bbZ$ is zero.

We have local points everywhere, and integral points at almost all
places by Lang's theorem \cite[Theorem~1 and its corollary]{Lang1956}
and smooth lifting. Borovoi's
Hasse-principle theorem for homogeneous spaces with connected
stabilizer therefore applies: its obstruction is the pairing with
the locally trivial algebraic Brauer group just shown to vanish
\cite[Theorem~2.2]{Borovoi}.  Thus $X(k)\ne\varnothing$, as required.
\end{proof}

Fix this $s$ and put
\[
 J_0=k(j),\qquad j=is,\qquad C_s=C_D(s).
\]
The field $LS_0$ is biquadratic over~$k$, with third quadratic field
$J_0$.  The algebra $C_s$ has degree~$m$ over~$LS_0$, and its
involution is unitary over~$S_0$.  We shall factor a representative~$u$
into a quaternion norm-one element and an element of~$\U(C_s)$.

\subsection{The factorization and its local norm conditions}

For $u\in S$, define
\begin{equation}\label{eq:unitary-factor-data}
 \begin{aligned}
 s'&=usu^{-1},& t&=s's^{-1},&
 p&=(u+sus^{-1})/2,\\
 z&=(1+t)/2=up^*,&
 h&=zz^*=pp^*=(2+t+t^{-1})/4.
 \end{aligned}
\end{equation}
These are identities in~$D$.  For example, expanding $4pp^*$ gives
$2+usu^{-1}s^{-1}+sus^{-1}u^{-1}$.  We have $p\in C_s$, and both $s$
and $s'$ invert~$t$ by conjugation.  Hence $h$ commutes with both of
them.

Over a splitting field, use the two $s$-eigenspaces to write
\[
 u=\begin{pmatrix}A&B'\\ C&D'\end{pmatrix}.
\]
The two diagonal blocks of~$h$ have the same characteristic polynomial
$\Psi$.  On the open where the blocks are invertible, the first is
\begin{equation}\label{eq:unitary-compression}
 h_+=A(u^{-1})_{++}
     =(1-B'(D')^{-1}CA^{-1})^{-1}.
\end{equation}
The analogous expression for the second block is conjugate to this
one, since the two products of the invertible off-diagonal block
maps are conjugate.  The equality of characteristic polynomials then
holds everywhere by polynomial continuation.  Exchange of the
$s$-labels preserves $\Psi$, and $h=h^*$ preserves it under the
unitary involution; thus $\Psi\in k[T]$.

Set $d_0=\Nrd_{C_s}(p^*)$.  Then $d_0\in J_0$: conjugation of
$LS_0/S_0$ takes it to $\Nrd_{C_s}(p)$, and exchange of the two
$s$-labels does the same, by the complementary-minor identities for
$u$ and $\det u=1$.  Their product therefore fixes~$d_0$, which is the
condition that it belong to~$J_0$.

We will choose $u$ so that $F=k(h)$ is a field of degree~$m$ and
$h(1-h)\ne0$.  When $m\geq3$, we also require that its normal closure
have group~$S_m$ even after adjoining~$LS_0$.  Taking determinants of
$pp^*=h$ gives
\begin{equation}\label{eq:unitary-norm-compatibility}
 N_{F/k}(h)=N_{J_0/k}(d_0).
\end{equation}
In~$D$, the elements
\[
 j=is,\qquad b=i\bigl(s'-(2h-1)s\bigr)
\]
anticommute and satisfy
\begin{equation}\label{eq:unitary-quaternion}
 j^2=i^2r,\qquad b^2=4j^2h(1-h).
\end{equation}
Indeed $ss'+s's=2r(2h-1)$, from which both assertions follow by
expansion.  Since both squares are nonzero, the quaternion algebra
they define over~$F$ is central simple, and its homomorphism into~$D$
is injective.  Denote its image by~$Q$.  The involution acts on~$Q$ as
quaternion conjugation.  Moreover $t=(is')j^{-1}$, so $z\in Q$ and
its quaternion norm is~$h$.  Scalar extension gives
$LQ=C_D(F)$, a central simple algebra of degree two over~$LF$.
Consequently $\Res_{F/k}\SL_1(Q)$ maps into~$S$.

Suppose that we find $a\in J_0F$ satisfying
\begin{equation}\label{eq:unitary-simultaneous-norms}
 N_{J_0F/F}(a)=h,
 \qquad N_{J_0F/J_0}(a)=d_0\quad\text{if }m\geq3.
\end{equation}
The subfield $J_0F$ is contained in both $C_s$ and~$Q$, and $*$ acts
on it as the involution over~$F$.  We obtain
\begin{equation}\label{eq:unitary-factorization}
 u=wq,\qquad w=za^{-1}\in\SL_1(Q),\qquad
 q=a(p^*)^{-1}\in\U(C_s)\cap S.
\end{equation}
The last membership follows because $wq=u$ and $w$ is unitary of
ambient determinant one.  For $m\geq3$, the second norm condition
also gives $\Nrd_{C_s}(q)=1$, so $q\in\SU(C_s)$.  We now construct
$u,a$ with the local control needed to kill these factors.

Choose a finite set $B$ containing all archimedean and fixed
exceptional places, the places below the rank-zero places of
$\SU(C_s)$ over~$S_0$, and the anisotropic finite places of~$S$.
At these places ask that $u$ be close to~$1$.  Additional places
split in~$LS_0$ allow us to prescribe all needed Frobenius cycle
types of~$F$.  The block matrices
\begin{equation}\label{eq:unitary-near-identity-model}
 u=\begin{pmatrix}I&\varepsilon I\\
                  \varepsilon X&I+\varepsilon^2X\end{pmatrix},
 \qquad h_+=I+\varepsilon^2X
\end{equation}
have determinant one and are arbitrarily close to~$1$ for small
$\varepsilon$.  Taking $X$ with any prescribed separable residue
factorization gives the desired local unramified algebra for~$F$.
These are open conditions on~$u$.  Prescribing representatives of
every conjugacy class of~$S_m$ at places split in~$LS_0$ ensures that
the Galois subgroup over~$LS_0$ is~$S_m$: a proper subgroup misses a
conjugacy class, by the derangement theorem for its coset action.
For $m=2$ a single irreducible quadratic residue condition gives the
required disjointness.  Include these auxiliary places in~$B$.

Near~$1$, the norm equations have local solutions near~$1$.  Start
with the analytic square root $a_0=h^{1/2}\in F\otimes k_v$; its norm
from $J_0F$ to~$F$ is~$h$.  For $m\geq3$, the ratio
\[
 c_0=d_0/N_{J_0F/J_0}(a_0)
\]
lies in $J_0^1$, by Equation~\eqref{eq:unitary-norm-compatibility},
and is close to~$1$.  Choose its $m$-th root $c\in J_0^1$ close to~$1$
and replace $a_0$ by~$ca_0$.  This solves both norm equations.  All
reduced eigenvalues involved are ambient eigenvalues of matrices
near~$1$, so $a,w,q$ can be made uniformly small in the sense of
Lemma~\ref{lem:unitary-uniform-roots}.

\subsection{Controlling every other place}

It remains to ensure local solvability of
Equation~\eqref{eq:unitary-simultaneous-norms} and splitting of~$Q$
away from the finitely many places at which the factors are made
small.  To apply Proposition~\ref{prop:power-approximation}, we first
specify the allowed local conditions.

Let $\Omega\subset S$ be the fixed Zariski open on which $\Psi$ is
separable and $h(1-h)$ is invertible.  For each place outside~$B$,
let $\mathcal U_v$ be the union of $S(k_v)\setminus\Omega(k_v)$,
the locus in $\Omega(k_v)$ on which $Q$ is split and the required
norm equations are soluble, and a sufficiently small generic
identity opening $I_v\subset\Omega(k_v)$ on which the analytic-root
construction supplies norm solutions with both factors $w,q$ local
$e$-th powers.  The preceding analytic construction makes $I_v$
nonempty.  Thus a point of $I_v$ is allowed even when its quaternion
algebra is not split.

The integral and pole tests below will give the split-$Q$ condition
for points in~$\Omega$.  Including the complement of~$\Omega$ in
$\mathcal U_v$ lets these tests meet the universal integral-point
condition of Proposition~\ref{prop:power-approximation} without adding
a discriminant divisor to its codimension-two exclusion.  We will
ensure that the final rational representative lies in~$\Omega$ by
one of the auxiliary local openings already prescribed.  After the
tests, we specify which exceptional places use the small opening
$I_v$ and which use the split-$Q$ condition.

At integral reduction, consider the four block-determinant
hypersurfaces
\[
 \det A=0,\quad\det B'=0,\quad\det C=0,\quad\det D'=0.
\]
Exclude their pairwise intersections.  If all four determinants are
nonzero and $m\geq3$, also require that $\Psi$ have at least one
simple root over the algebraic closure of the residue field.  If
exactly one block determinant vanishes, require, for $m\geq3$, that
$0$ be a simple root of~$\Psi$ for a diagonal block, or $1$ a simple
root for an off-diagonal block.  The latter roots follow respectively
from the block and complementary-minor determinant identities for
$h$ and~$1-h$.

All excluded loci have geometric codimension at least two.  This can
be checked first on~$\GL_{2m}$: the four block determinants define
distinct irreducible hypersurfaces.  On the open where they are all
invertible, the map
$u\mapsto B'(D')^{-1}CA^{-1}$ is a submersion, as is seen by varying
$B'$ with the other three blocks fixed.  The transformation
$M\mapsto(1-M)^{-1}$ in Equation~\eqref{eq:unitary-compression}
preserves eigenvalue multiplicities on this open.  For $m\geq3$, the locus of
monic degree-$m$ polynomials with no simple root has dimension at most
$\lfloor m/2\rfloor$ in coefficient space and thus codimension at
least two.  The characteristic-polynomial map on matrices is flat:
each fiber has dimension $m^2-m$, by the decomposition into finitely
many Jordan types and their centralizer dimensions.  Thus the same
codimension bound holds for its inverse image.  On each single block
hypersurface the requisite simple root holds on a nonempty open, as
one sees from independent two-by-two transformations between paired
eigenlines, with only one of the relevant scalar blocks vanishing.
Finally, scalar multiplication changes none of the zero or root
conditions, so passage from $\GL_{2m}$ to $\SL_{2m}$ preserves these
codimensions.  The union is Galois stable and spreads to the chosen
integral models after finitely many exceptions.

Consider first an integral place outside the fixed exceptions where
$J_0$ is inert.  Frobenius interchanges the two diagonal block
divisors and also the two off-diagonal block divisors.  To check
this, exchange of the $s$-labels swaps each pair, and the unitary
diagram action does the same: in adapted dual bases it is inverse
transpose, and complementary minors swap the paired determinants.
Exactly one of these actions occurs when the third quadratic field
$J_0$ is inert.  A rational residue point therefore cannot lie on
exactly one divisor.  All four determinants are nonzero.  Hence $h$
and $1-h$ are units in every local component of~$F$.  The quadratic
algebra $J_0F/F$ is unramified or split, so these units are norms.
In particular Equation~\eqref{eq:unitary-quaternion} shows that $Q$
splits, and one can choose unit preimages of~$h$ for the first norm
equation.

For $m\geq3$ the simple residue root of~$\Psi$ supplies an
unramified local factor $E/k_v$ of~$F$.  It need not be a degree-one
factor.  Write $K=k_v$ and $J=J_0\otimes_k K$, an unramified
quadratic field in the case under consideration.  If $a_0$ denotes
the chosen unit preimages of~$h$, the required correction is
\[
 \eta=d_0/N_{JF/J}(a_0)\in\mathcal O_J^1,
\]
by Equation~\eqref{eq:unitary-norm-compatibility}.  Hilbert~90 writes
$\eta=y/\bar y$ with $y\in J^\times$.  Since $J/K$ is unramified,
multiplication by a power of a $K$-uniformizer makes $y$ a unit
without changing this quotient.  The unramified etale algebra
$JE/J$ has surjective norm on units.  Choose a unit $x\in JE$ with
$N_{JE/J}(x)=y$, and put $\beta=x/x^*$.  Then
\[
 N_{JE/E}(\beta)=1,\qquad N_{JE/J}(\beta)=\eta.
\]
Extend $\beta$ by~$1$ on the other factors of $JF$ and replace
$a_0$ by $a_0\beta$.  This corrects the determinant while preserving
the norm to~$F$, for either parity of $[E:K]$.

If $J_0$ splits, then $Q$ splits automatically.  Write
$a=(a_+,a_-)$ and $d_0=(d_+,d_-)$.  Choose
$N_{F/K}(a_+)=d_+$ and set $a_-=h/a_+$; compatibility gives
$N_{F/K}(a_-)=d_-$.  Thus one component must have a prescribed norm
from~$F$.
When all block determinants are nonzero, the prescribed value is a
unit, so an unramified factor supplies it by ordinary unit norm
surjectivity.  If one determinant vanishes, the prescribed simple
residue root is literally $0$ or~$1$, so Hensel's lemma supplies a
degree-one factor of~$F$.  Its norm map supplies any prescribed value,
including a nonunit.  This proves all the local assertions at the
permitted integral reductions.  When $m=2$, no determinant norm is
being imposed, so the splitting and first-norm assertions already
suffice.

We must check the single poles permitted by
Proposition~\ref{prop:power-approximation}.  Enlarge its fixed
splitting field to split $L,S_0,D$ and all the fixed data.  At a
single pole the place splits in this field, and
\[
 u=g_L\diag(t_{\rm p}^{ed_1},\ldots,t_{\rm p}^{ed_{2m}})g_R,
 \qquad d_1<\cdots<d_{2m},\quad \ord_v(t_{\rm p})>0,
\]
with integral side factors.  For $m\geq3$ impose total splitting of
$\Psi$ by requiring finitely many nonzero leading minors.  Explicitly,
if $c_b$ is its $b$-th elementary coefficient, exterior-power
expansion of the compression in
Equation~\eqref{eq:unitary-compression} gives
\begin{equation}\label{eq:unitary-pole-valuations}
 \ord_v(c_b)=e\ord_v(t_{\rm p})
       \sum_{a=1}^b(d_a-d_{2m+1-a}),\qquad 1\leq b\leq m.
\end{equation}
For completeness, the least exponent uses the first $b$ exponent
indices in the positive exterior power and the last $b$ in its
inverse.  Its coefficient is the product of the corresponding
opposite minors in the two rank-$m$ projections
\[
 g_RP_+g_R^{-1},\qquad g_L^{-1}P_+g_L,
\]
where $P_+$ is projection onto one $s$-eigenspace.  All other terms
have larger valuation.  The indicated minors are nonzero for a
generic conjugate of a rank-$m$ projection; require them all to be
units on the pole hyperplane.  These are nonempty simultaneous open
conditions.  Their choices can be made in the staged order required
by Proposition~\ref{prop:power-approximation}.  With all other torus
parameters equal to~$1$, the boundary of the factor with conjugator
$k_j$ has $g_L=xk_j$ and $g_R=k_j^{-1}$.  First choose the basic
conjugators so that every right-projection minor for $k_j^{-1}$ is
nonzero, along with the basic-list dominance conditions.  Then choose
$x$ in the intersection of the nonempty left-minor opens for the
finitely many $xk_j$.  All coordinate conjugates give only finitely
many such dense open conditions.  For an appended factor, with $x$
already fixed, its conjugator $k$ can be chosen in the intersection
of the right-minor open and the inverse translate by $x$ of the
left-minor open.  This intersection is again dense open and meets
the required local approximation openings.  Setting each new
parameter to~$1$ preserves all the old boundary tests.  This proves
nonemptiness of every single-pole test without changing the earlier
choices.

The consecutive slopes in Equation~\eqref{eq:unitary-pole-valuations}
are strictly increasing.  Every Newton-polygon segment has horizontal
length one, so $\Psi$ splits completely over~$k_v$.  Since $J_0$
also splits here, all required norm equations are soluble and~$Q$
splits.  For $m=2$, splitting of~$J_0$ alone suffices.

The tests establish the required membership in~$\mathcal U_v$ at
permitted integral reductions and single poles.  It remains to
retain global genericity and to choose the openings at the finite
exceptional places.  The final rational representative lies
in~$\Omega$ because one of the existing auxiliary local openings is
contained in~$\Omega(k_v)$.  In
Equation~\eqref{eq:unitary-near-identity-model}, choose $X$ with
separable characteristic polynomial and with both $X$ and
$I+\varepsilon^2X$ invertible.  Then $h_+=I+\varepsilon^2X$ has
distinct eigenvalues different from $0,1$, and these conditions
persist in its local opening.  This argument uses the local matrix
itself, not the reduction of~$\Psi$, which can have repeated roots
when $\varepsilon$ is small.
After the initial representative $x$ is chosen, let $B_1$ be its
finite set of denominator places outside~$B$.  At these places
specifically choose the opening $I_v\subset\mathcal U_v$.
This requires no splitting assumption; the block model in
Equation~\eqref{eq:unitary-near-identity-model} was needed only to
prescribe Frobenius types at the auxiliary split places.  The finitely
many later model exceptions are assigned unit-parameter
reductions with all four block determinants nonzero and separable
$\Psi$; such opens are nonempty over sufficiently large residue
fields, with small fields already included among the fixed
exceptions, and thus use the split-$Q$ branch as well.  Consequently
the small branch is used only at the controlled finite set
$B^+=B\cup B_1$.  The construction produces a representative of
the required coset for which $Q$ is split outside~$B^+$.
At every place Equation~\eqref{eq:unitary-simultaneous-norms} is
soluble with the prescribed near-identity solutions at~$B^+$.

For $m\geq3$, the simultaneous equations form a torsor under the
norm torus of Lemma~\ref{lem:norm-torus}, with quadratic field~$J_0$
and degree-$m$ field~$F$.  The full symmetric action over~$LS_0$
implies the joint $S_m\times C_2$ hypothesis for $F,J_0$.  That lemma
gives a global~$a$ with the local approximations.  When $m=2$, the
quadratic Hasse norm theorem over~$F$
\cite[Chapter~VIII, Theorem~3.1]{MilneCFT} and weak approximation on its
norm torsor give~$a$.  In either case the quaternion factor~$w$ in
Equation~\eqref{eq:unitary-factorization} is killed: all its
rank-zero places lie over the controlled set, where it is a local
$e$-th power by Lemma~\ref{lem:unitary-uniform-roots}.  If $m\geq3$,
the factor $q$ lies in the smaller group $\SU(C_s)$ over~$S_0$ and is
small at every rank-zero place of that fixed group.  Induction and
Lemma~\ref{lem:restriction-mp} kill~$q$ as well.

\subsection{The determinant correction in degree four}

It remains to treat $m=2$.  The preceding construction has killed~$w$
and produced $q\in\U(C_s)\cap S$, small at the fixed set~$B$, but its
determinant in~$C_s$ may not be~$1$.  Write
\[
 r_0=\Nrd_{C_s}(q).
\]
Both conjugations of $LS_0$ over $S_0$ and over~$L$ invert~$r_0$:
the first by unitarity, the second by ambient determinant one.  Thus
$r_0\in J_0^1$.  Taking $q$ sufficiently close to~$1$ at~$B$ ensures
$r_0\ne-1$.  Put
\[
 d=(1+r_0)/2\in J_0^\times,\qquad r_0=d/\bar d,
\]
where the bar is the nontrivial automorphism of~$J_0/k$.  The element
$d$ is close to~$1$ at~$B$.

We will remove this determinant by two elements in smaller binary
special groups.  The following construction gives the freedom needed
to keep their varying rank-zero places under control.  For any finite
set of places outside the fixed exceptions, one can choose a
quadratic field $F_i=k(f_i)\subset C_s$, disjoint from~$LS_0$, with
$f_i^*=f_i$, such that $F_i$ splits at these places and the binary
group $\SU(C_D(F_i),*)$ is split over both local components of~$F_i$.
Here and below $B$ includes all fixed exceptions of the structures
defined from~$s$.

To prove the construction, let $B_s$ be the quaternion algebra over
$S_0$ obtained by descent from~$C_s$: its defining semilinear
involution is the composition of $*$ and standard quaternion
conjugation.  These involutions commute, and
$B_s\otimes_{S_0}LS_0=C_s$.  For a trace-zero element $y\in B_s$,
the element $f=i^{-1}y$ satisfies $f^*=f$.  Prescribe
$f^2=c\in k^\times$, so $y^2=i^2c$.  A quadratic etale algebra
embeds into $B_s$ if and only if it is a field at every ramified
place of~$B_s$, by the local embedding criterion and Brauer
invariants.  We may impose this by choosing $i^2c$ nonsquare there.
At a finite place of~$k$ the kernel of the square-class map to a
quadratic component of~$S_0$ has at most two elements; in the split
case it has one.  Thus the forbidden square classes do not exhaust
the local square classes.  At real places the requirement is simply
a sign.  These ramified places are above~$B$.  At the prescribed
places outside~$B$ require $c$ to be a square, and at one further
place split completely in~$LS_0$ require it to be a nonsquare.
Weak approximation supplies~$c$.  The last condition ensures that
$F_i$ is a field disjoint from~$LS_0$.

The embedding of $S_0(\sqrt{i^2c})$ into~$B_s$ can approximate any
prescribed local embeddings: they are conjugate, and weak
approximation in the open affine variety $B_s^\times$ approximates
the conjugators.  At a specified place split in~$L$, ordinary split
block embeddings make the two binary centralizers split.  At a good
unramified nonsplit place where $S_0$ splits, decompose both
two-dimensional $s$-blocks into orthogonal lines with unit norm
values.  Give each $f_i$-eigenspace one line from each block.  Its
Hermitian determinant is a unit, and hence a norm, so each binary
space is hyperbolic.  If $S_0$ does not split there, it locally equals
$L$; the two $s$-blocks are dual totally isotropic spaces.  Choose
paired line decompositions and assign one hyperbolic plane to each
$f_i$-eigenvalue.  In each case $f_i$ is Hermitian and traceless in
$C_s$, and the desired local property persists on an open set of
embeddings.  This proves the construction.

For such a field $F_i$, the compositum $LS_0F_i$ is a maximal
subfield of~$C_s$, and there is an inclusion
\begin{equation}\label{eq:unitary-binary-norm-tori}
 (J_0F_i/F_i)^1
 \ \subset\ \SU(C_D(F_i),*)\cap\U(C_s).
\end{equation}
Indeed the determinant in the binary algebra $C_D(F_i)$ is the norm
from $LS_0F_i$ to~$LF_i$.  On $J_0F_i$ its nontrivial automorphism
exchanges the $s$-labels, which is the norm-one condition displayed
on the left.  The determinant of the same element in~$C_s$ is its
norm from $J_0F_i$ to~$J_0$.

Choose $F_1$ first.  Outside~$B$, all but finitely many places have
both of the following properties: $d$ is a local norm from
$J_0F_1$ to~$J_0$ at every place above the given $k$-place, and the
binary group in Equation~\eqref{eq:unitary-binary-norm-tori} is
isotropic at every place of~$F_1$ above it.  This follows from
unramified unit norm surjectivity and good reduction.  Call these
the places covered by~$F_1$.  Choose $F_2$ by the preceding
construction to split, with split binary groups in both components,
at every remaining place outside~$B$.  Every such remaining place
is covered by~$F_2$, since the local norm from a split quadratic
algebra is surjective.

Solve
\begin{equation}\label{eq:unitary-multinorm-correction}
 N_{J_0F_1/J_0}(x_1)N_{J_0F_2/J_0}(x_2)=d.
\end{equation}
This equation satisfies the Hasse principle and weak approximation
on its locus where the variables are invertible.  Indeed putting
$y_2=x_2^{-1}$ turns it into the nonzero locus of the nonsingular
four-variable quadratic equation
\[
 N_{J_0F_1/J_0}(x_1)-dN_{J_0F_2/J_0}(y_2)=0.
\]
Hasse--Minkowski \cite[Chapter~VIII, Theorem~3.5(b)]{MilneCFT}
gives a rational isotropic vector whenever there
are local ones.  The corresponding projective quadric is rational
once it has a point, and its affine cone minus the vertex has weak
approximation.  Removing the two norm-zero loci preserves the
needed approximation to local points in this open set.

The requisite local points can be chosen with additional control.
At~$B$ take $x_1=\sqrt d$ and $x_2=1$, near~$1$; the square root is
in the local $J_0$-algebra.  Outside~$B$, use one covering field for
the entire $k$-place to carry the norm~$d$, and take the other
variable to be~$1$ at every component above that place.  In
particular, at any rank-zero place of one of the two binary groups,
its own variable can be taken to be~$1$: that field cannot be the
covering field there.  These rank-zero places form a finite set.
Use weak approximation in
Equation~\eqref{eq:unitary-multinorm-correction} at their underlying
places and at~$B$, obtaining global invertible~$x_1,x_2$ with these
local approximations.

Set $b_i=x_i/(x_i^*)$.  By
Equation~\eqref{eq:unitary-binary-norm-tori}, each $b_i$ belongs to
the indicated binary special group, of reduced degree two over
$F_i$, and
\[
 \Nrd_{C_s}(b_1)\Nrd_{C_s}(b_2)
   =\frac{d}{\bar d}=r_0.
\]
At every rank-zero place of its binary group, $b_i$ is as close to~$1$
as necessary.  The $A_1$ case, local power lifting, and
Lemma~\ref{lem:restriction-mp} therefore kill~$b_1,b_2$.  The element
$(b_1b_2)^{-1}q$ has determinant one in~$C_s$ and remains close to~$1$
at~$B$.  Its special group over~$S_0$ has degree two and all its
rank-zero places lie above~$B$, so it too is killed.  Together with
the already killed factor~$w$, this kills~$u$.

The even induction step is complete.  The degree-two base, the
odd-degree division result, the odd matrix argument, and this step
prove $V_0/R=1$ in every reduced degree over every number field.
Proposition~\ref{prop:finite-defect} now proves
Theorem~\ref{thm:outer-a}.
